% Шаблон документа требований: «Системные требования», «Требования к ПО».
% Скопируйте в docs/0N_<имя>.tex, замените название, версию и дату, заполните разделы.
% Красные пометки \clarify{…} — места, которые надо заполнить; перед сдачей их не остаётся.
% Правила оформления — docs/00_formatting.tex. Собрать: ./build.sh docs/0N_<имя>.tex
\documentclass{tpks}
\doctitle{Название документа}
\docversion{1}
\docdate{ДД.ММ.ГГГГ}

\begin{document}
\makeheader
% Лист регистрации изменений — сразу после \makeheader. Новая версия — новая строка.
\begin{changelog}
\change{1}{ДД.ММ.ГГГГ}{Первая версия}{Фамилия И. О.}{Задание к паре ДД.ММ.ГГГГ}
\end{changelog}

\section{Введение}
\subsection{Назначение}
\clarify{что определяет документ и кто им пользуется}

\subsection{Область применения}
\clarify{границы системы или ПО, к которым относятся требования}

\subsection{Термины и сокращения}
Определения взяты из общего глоссария проекта.
\printterms{ПО}

\subsection{Ссылки}
\begin{enumerate}
  \item \clarify{«Название документа». Что из него используется}
\end{enumerate}

\subsection{Обзор}
\clarify{что содержат следующие разделы}

\section{Общее описание}
\subsection{Окружение}
\clarify{схема «система — пользователи — внешние системы»; внешние системы названы по функции}

\subsection{Пользователи}
\clarify{роли пользователей}

\subsection{Ограничения, допущения и зависимости}
\clarify{ограничения с обоснованием}

\section{Требования}
% Порядок разделов — README.md §4.1 (системные требования) или §4.5 (требования к ПО).
\subsection{Функциональные требования}

\begin{req}{SYS_01_01}{1}
Система должна \clarify{что делает система, а не как}.
\reqsource{\clarify{потребность, вариант использования или системное требование}}
\reqverify{\clarify{испытание, анализ или рассмотрение и критерий «пройдено / не пройдено»}}
\end{req}

\subsection{Нефункциональные требования}
\clarify{время отклика, нагрузка, надёжность — в числах}

\section{Приложение: иллюстрации}
\clarify{сценарии и схемы, которые не являются требованиями; раздел можно удалить}

\end{document}
